\tikzsetfigurename{00_preface-}
%% 00_preface.tex -- Preface (front matter, unnumbered)
\chapter{Preface}\label{ch:preface}

A heartbeat on the skin is about a millivolt. A brain rhythm on the scalp is a few tens of microvolts. A single neuron, seen from an electrode a few tens of micrometres away, is a hundred microvolts that lasts a millisecond. Every wearable ECG patch, EEG headset, prosthetic-control armband and neural implant has to turn these signals into clean numbers while the patient's body picks up volts of mains interference, the electrodes add hundreds of millivolts of offset, and the battery allows a few microwatts per channel. The chip that does this is a \textbf{biopotential analog front-end} (AFE) and its data converters.

This book teaches you to design one by following a single, concrete chip, \textbf{\OB{}}: an 8-channel ECG/EEG/EMG front-end with a 12-bit SAR and a 16-bit incremental $\Delta\Sigma$ converter, built in the open-source IHP \sgx{} \qty{130}{\nano\meter} SiGe BiCMOS process~\cite{ihppdk} with open-source tools~\cite{iicosic}. Every concept follows the same pattern as our \emph{Open-SERDES} tutorial~\cite{serdestut}: a short derivation, a worked example with the project's real numbers, and a \emph{what this means for your design} box.

\paragraph{Who this book is for} Third- and fourth-year electrical engineering undergraduates and first-year graduate students, and biomedical engineers who need to specify the hardware they will use. After reading it you should be able to:
\begin{itemize}
  \item turn a clinical or research application into a front-end specification: channels, bandwidth, noise, CMRR, input impedance, offset range and power;
  \item design and budget a low-noise, low-power amplifier chain: chopping, DC servo loops, filters and right-leg drive;
  \item choose, size and verify a SAR or $\Delta\Sigma$ converter for a biosignal;
  \item show, with real ECG, EEG and EMG recordings, how the circuit choices move a clinical result;
  \item carry a block from schematic to a verified layout with open tools, ready for an IHP OpenMPW shuttle.
\end{itemize}

\paragraph{Prerequisites} Circuit analysis, signals and systems (Fourier and Laplace transforms, filters) and one course in analog electronics (small-signal MOS and BJT models, the differential pair, op-amps), plus some Python. The transistor-level background of Part~II, in particular the $\gm/I_\mathrm{D}$ method, is taught in Part~II of \emph{The Art of Analog IC Design}~\cite{artofaicd}; we cross-link it instead of repeating it.

\section*{How the book is organised}
The book has five parts. \textbf{Part~I, Foundations} (Chapters~1--4), explains where biopotentials come from, what the electrodes and the body do to them, and the interference, safety rules and standards every front-end must meet. \textbf{Part~II, The Analog Front-End} (Chapters~5--12), designs the low-noise, low-power amplifier chain. \textbf{Part~III, Data Converters} (Chapters~13--22), is the heart of the book: it takes you from the quantiser to a measured SAR and $\Delta\Sigma$ converter. \textbf{Part~IV, Biosignal Processing and Systems} (Chapters~23--28), processes real ECG, EEG and EMG data, puts processing on chip and adds bioimpedance sensing. \textbf{Part~V, Putting It Together} (Chapters~29--31), assembles the architecture and its budgets, walks through the open-source flow and collects the fifteen labs.

\paragraph{Three reading paths} Everyone should read Part~I first. After that, \cref{fig:pre-paths} suggests three routes. The \textbf{circuit path} (Parts I, II, III and V) is for students heading to a tape-out. The \textbf{signal path} (Part~I, Part~IV, the converter specifications of \cref{ch:sampling,ch:adcspecs}, then Part~V) is for biomedical-engineering and DSP students who need to specify, not build, the hardware. The \textbf{converter-only path} (\cref{ch:sampling} to \cref{ch:adctest} with Labs~6--10) is a self-contained course on data converters.

\begin{figure}[htbp]
\centering
\tikzset{pnode/.style={draw=#1,fill=white,line width=0.8pt,rounded corners=2pt,minimum height=7mm,minimum width=13mm,
    inner sep=2pt,font=\sffamily\footnotesize\bfseries,text=#1,align=center},
  plink/.style={line width=1.1pt,draw=#1}}
\begin{tikzpicture}[x=1cm,y=1cm]
\node[pnode=navy] (p1) at (0,0) {Part I\\[-2pt]\scriptsize Ch 1--4};
% circuit path
\node[pnode=golddk] (c2) at (2.6,1.3) {Part II\\[-2pt]\scriptsize Ch 5--12};
\node[pnode=golddk] (c3) at (5.0,1.3) {Part III\\[-2pt]\scriptsize Ch 13--22};
% signal path
\node[pnode=bblue] (s4) at (2.6,0) {Part IV\\[-2pt]\scriptsize Ch 23--28};
\node[pnode=bblue] (s13) at (5.0,0) {Ch 13--14\\[-2pt]\scriptsize ADC specs};
% converter-only path
\node[pnode=slate] (k1) at (2.6,-1.3) {Ch 13--22};
\node[pnode=slate] (k2) at (5.0,-1.3) {Labs 6--10};
% merge
\node[pnode=navy] (p5) at (7.6,0.65) {Part V\\[-2pt]\scriptsize Ch 29--31};
\draw[plink=golddk] (p1) -- (c2); \draw[plink=golddk] (c2) -- (c3); \draw[plink=golddk] (c3) -- (p5);
\draw[plink=bblue] (p1) -- (s4); \draw[plink=bblue] (s4) -- (s13); \draw[plink=bblue] (s13) -- (p5);
\draw[plink=slate] (k1) -- (k2);
\node[font=\sffamily\scriptsize\bfseries,text=golddk,anchor=south] at (3.8,1.7) {circuit path: to a tape-out};
\node[font=\sffamily\scriptsize\bfseries,text=bblue,anchor=south] at (3.8,0.36) {signal path: specify the hardware};
\node[font=\sffamily\scriptsize\bfseries,text=slate,anchor=north] at (3.8,-1.72) {converter-only path: a data-converter course};
\end{tikzpicture}
\caption{\textbf{Three ways to read the book.} The circuit path (gold) follows the front-end and converter circuits to a tape-out. The signal path (blue) goes from the signals to their processing and the converter specifications. The converter-only path (grey) is Part~III with its labs on its own.}
\label{fig:pre-paths}
\end{figure}

\section*{How to read the boxes}
Chapters open with what you will learn and close with key takeaways. In between, six kinds of boxes mark the role of each piece of text:

{\renewcommand{\thewexample}{P.1}%
\begin{keyeq}[Signal-to-noise ratio in decibels]
\[ \SNR = 20\log_{10}\frac{V_\mathrm{signal,rms}}{V_\mathrm{noise,rms}} \]
\end{keyeq}
\begin{example}{A number you will use on every page}
A \qty{20}{\micro\volt} rms EEG rhythm over the \qty{1.5}{\uVrms} noise target of \OB{} gives $\SNR = 20\log_{10}(20/1.5) \approx \qty{22.5}{\deci\bel}$.
\end{example}}
\setcounter{wexample}{0}
\begin{insight}
An \emph{insight} box says what a result means for the \OB{} design.
\end{insight}
\begin{pitfall}
A \emph{pitfall} box warns about a common mistake, such as quoting an amplitude in rms when the specification is peak-to-peak.
\end{pitfall}
\begin{tryit}[Lab 1: biosignal explorer]
A \emph{try-it} box points to a lab script (\cref{sec:lab1}).
\end{tryit}
\begin{specnote}
A \emph{specnote} box states an \OB{} target: input-referred noise below \qty{1.5}{\uVrms} from \qtyrange{0.5}{100}{\hertz} in EEG mode.
\end{specnote}
Longer algebra sits in grey \emph{derivation} boxes.

\section*{How the labs work}
Fifteen labs in three tracks use only free tools: Python (NumPy, SciPy, \tool{wfdb}, \tool{python-deltasigma}) for signals and system models, \tool{Xschem} and \tool{ngspice} for transistor-level \sgx{} circuits, and \tool{KLayout} for layout (\cref{ch:labs}). The Python labs live in \code{labs/python/} and share one module, \code{biosig.py}, which holds the signal tables, the synthetic signals and the electrode and interference models. The figure scripts in \code{tutorial/figsrc/} import the same module, so the numbers in the book and the numbers your labs print come from the same code. Datasets stay on PhysioNet~\cite{physionet} and NinaPro~\cite{ninapro} under their own licences; the labs download them on first use.

\section*{The state of this draft}
The book is being written in six stages. This first release (v0.1) contains the complete outline of all 31 chapters, the four chapters of Part~I and Labs~1 and~2 in full, and the appendices. Chapters still in outline form say so at the top and list their topics, their design question and their labs. The next stage adds the Python-only converter core (\cref{ch:sampling,ch:adcspecs,ch:dstheory} with Labs~6 and~9).

\section*{Sources and acknowledgments}
The biomedical-instrumentation chapters draw on Webster's \emph{Medical Instrumentation}~\cite{webster} and on Yazicioglu, Van Hoof and Puers' \emph{Biopotential Readout Circuits}~\cite{yazicioglubook}; the converter chapters on Pelgrom~\cite{pelgrom} and on Schreier, Temes and Pavan~\cite{schreier}. The chip is possible because IHP opened its \sgx{} process design kit~\cite{ihppdk} and the Institute for Integrated Circuits at JKU Linz maintains the IIC-OSIC-TOOLS container and the analog chip template~\cite{iicosic,jkutemplate}. The labs use open data from PhysioNet~\cite{physionet} and NinaPro~\cite{ninapro}.

This is a \textbf{living draft (v0.1)}; numbers will be refined as the design matures. If you find an error, an unclear derivation or a figure that does not explain what it should, please email \href{mailto:bricd.org@gmail.com}{bricd.org@gmail.com} or open an issue or a pull request on the project's GitHub repository, linked from thebricd.org.

\section*{How this book was made}
This book, its figures and its lab code were developed with the help of AI tools, and this draft may contain mistakes; please report any you find as described above, and we will correct them. The topics, structure and many of the results follow the sources acknowledged above and listed in the References. Standards are summarised, never reproduced: always work from the edition of IEC~60601 that applies to your device.
